% =============================================================================
% 2. fejezet — Szakirodalom és elméleti háttér
% -----------------------------------------------------------------------------
% A) SZAKIRODALOM (ez a rész)
% B) ELMÉLETI HÁTTÉR — a következő lépésben kerül a fájl végére
% =============================================================================

Ez a fejezet két nagyobb egységből áll. Az első rész a bőrelváltozások
gépi tanulással történő osztályozásának szakirodalmi előzményeit tekinti át --
a klinikai diagnosztikai gyakorlattól a klasszikus gépi tanuláson és a
konvolúciós neurális hálókon át a Vision Transformer és a multimodális nagy
nyelvi modellek megjelenéséig. A második rész a dolgozatban felhasznált
kulcstechnológiák elméleti hátterét ismerteti.

% =============================================================================
\section{A bőrrák és a korai diagnózis jelentősége}
% =============================================================================

A bőrrák a leggyakrabban diagnosztizált rosszindulatú daganatok közé tartozik
világszerte. A Nemzetközi Rákkutató Ügynökség (IARC) GLOBOCAN-becslései szerint
2022-ben több mint 1,5 millió új bőrrákos esetet regisztráltak, ezen belül a
melanomák száma megközelítőleg 331~700 volt, mintegy 58~700 halálesettel
\cite{iarc2022melanoma}. Bár a melanoma a bőrrákos esetek kisebb hányadát teszi
ki, agresszív viselkedése és áttétképző hajlama miatt a bőrrák okozta halálozás
túlnyomó részéért felel.

A betegség kimenetele döntően a felismerés stádiumától függ. A még lokalizált
stádiumban diagnosztizált melanoma ötéves relatív túlélési aránya közel 100\%,
míg távoli áttét megjelenése esetén ez mintegy 35\%-ra csökken
\cite{seer2024melanoma}. Ez a markáns különbség teszi a korai felismerést
népegészségügyi szempontból kiemelten fontossá: a diagnózis néhány hónapos
késedelme is jelentősen ronthatja a beteg túlélési esélyeit.

\subsection{A klinikai diagnosztikai folyamat}

A melanoma klinikai gyanúja jellemzően vizuális vizsgálat alapján merül fel,
amelyet dermatoszkópos értékelés egészít ki; a végleges diagnózist minden
esetben szövettani vizsgálat (hisztopatológia) erősíti meg \cite{garbe2022melanoma}.
A magas kockázatú betegek nyomon követésében a szekvenciális digitális
dermatoszkópia és a teljes testfelszín fényképezése segíti a korai melanomák
kiszűrését \cite{garbe2022melanoma}. Ez a folyamat magas szintű szakmai
tapasztalatot és speciális eszközöket igényel.

A pigmentált bőrelváltozások szabad szemmel történő, gyors megítélését a
klinikai gyakorlatban az úgynevezett ABCDE-szabály támogatja. A betűszó a
melanomára jellemző öt figyelmeztető jegyet foglalja össze: aszimmetria
(\textit{Asymmetry}), szabálytalan határvonal (\textit{Border}), színbeli
heterogenitás (\textit{Color}), 6~mm-t meghaladó átmérő (\textit{Diameter}),
valamint az elváltozás időbeli változása (\textit{Evolving}); az utóbbi, ötödik
kritériummal Abbasi és munkatársai egészítették ki az eredeti ABCD-szabályt,
hangsúlyozva a fejlődő elváltozások jelentőségét \cite{abbasi2004abcde}. Az
ABCDE-szabály egyszerűsége miatt mind a szakemberek, mind a betegek számára jól
használható, ugyanakkor önmagában nem helyettesíti a műszeres és szövettani
vizsgálatot.

% MEGJEGYZÉS (Ervin): ide jól illene egy ábra az ABCDE-szabály szemléltetésére
% (pl. benignus vs. malignus elváltozások összehasonlítása). Ha van megfelelő,
% jogtiszta képed, cseréld le az alábbi placeholder dobozt.
\begin{figure}[H]
  \centering
  \fbox{\begin{minipage}[c][4cm][c]{0.85\textwidth}
    \centering\itshape
    [Ábra helye] -- Az ABCDE-szabály szemléltetése: jóindulatú és rosszindulatú
    pigmentált elváltozások összehasonlítása. (Saját vagy jogtiszta forrásból
    származó ábrát illessz be.)
  \end{minipage}}
  \caption{Az ABCDE-szabály vizuális kritériumai.}
  \label{fig:abcde}
\end{figure}

\subsection{A hozzáférés korlátai}

A bőrgyógyászati szakellátáshoz való hozzáférés világszerte egyenetlen.
Vidéki térségekben és alacsonyabb jövedelmű országokban a szakorvosi
kapacitás korlátozott, ami késedelmes diagnózishoz vezet; de még fejlett
egészségügyi rendszerekben is jelentős várakozási idővel kell számolni a
szakvizsgálatig. Ez a hozzáférési szakadék motiválja az olyan automatizált,
akár okostelefonon is futó szűrőeszközök fejlesztését, amelyek elsődleges
triázsként jelezhetik, ha egy elváltozás szakorvosi vizsgálatot indokol.

% =============================================================================
\section{Klasszikus gépi tanulás a bőrelváltozások osztályozásában}
% =============================================================================

A mélytanulás elterjedése előtt a bőrelváltozás-osztályozó rendszerek jellemzően
kétlépcsős felépítést követtek: először kézzel tervezett jellemzőket
(\textit{handcrafted features}) nyertek ki a képből, majd ezeken egy klasszikus
gépi tanulási osztályozót -- például tartóvektor-gépet (SVM), döntési fát vagy
véletlen erdőt (Random Forest) -- tanítottak be. A jellemzők gyakran közvetlenül
az orvosi szempontokat tükrözték: az ABCDE-szabályhoz kapcsolódó alaki
(aszimmetria, határvonal), színbeli (szín-momentumok, hisztogramok) és textúra
jellemzőket (például a szürkeszintű együttes előfordulási mátrixból, GLCM
számított mutatókat).

Kassem és munkatársai szisztematikus áttekintése a terület módszertani
fejlődését dokumentálja, és rámutat, hogy a klasszikus, kézi jellemzőkön alapuló
megközelítések teljesítményét erősen korlátozza a jellemzőtervezés minősége és a
felvételi körülmények (megvilágítás, felbontás, műtermékek) iránti
érzékenység \cite{kassem2021review}. Ezek a rendszerek jellemzően jól
kontrollált, dermatoszkópos felvételeken működtek megbízhatóan, általánosító
képességük azonban korlátozott maradt. Mindazonáltal a kézi jellemzőkön alapuló
osztályozók a mai napig értékes \textit{viszonyítási alapként} (baseline)
szolgálnak, mivel gyorsak, értelmezhetők, és megmutatják, mekkora teljesítmény
érhető el a reprezentáció automatikus tanulása nélkül.

% =============================================================================
\section{A mély tanulás megjelenése: konvolúciós neurális hálók}
% =============================================================================

A konvolúciós neurális hálók (CNN-ek) áttörést hoztak a képosztályozásban azzal,
hogy a jellemzőkinyerést és az osztályozást egyetlen, végponttól végpontig
tanítható modellbe integrálták, megszüntetve a kézi jellemzőtervezés
szükségességét. A bőrgyógyászati alkalmazások egyik úttörő munkájában Esteva és
munkatársai kimutatták, hogy egy nagy képadatbázison előtanított CNN a bőrrák
osztályozásában eléri a tapasztalt dermatológusok teljesítményét
\cite{esteva2017dermatologist}. Ezt követően több vizsgálat is megerősítette,
hogy a mélytanulás-alapú rendszerek megközelíthetik, sőt bizonyos, jól
körülhatárolt feladatokban meg is haladhatják az emberi szakértők
teljesítményét: Haenssle és munkatársai egy CNN-t 58 dermatológussal vetettek
össze dermatoszkópos melanoma-felismerésben \cite{haenssle2018man}, Brinker és
munkatársai pedig arról számoltak be, hogy mély neurális hálók a melanoma
képi osztályozásában felülmúlták a bevont bőrgyógyászokat \cite{brinker2019deep}.

E rendszerek teljesítményének kulcsa az \textit{átviteli tanulás} (transfer
learning): a modelleket először egy nagy, általános képadatbázison -- jellemzően
az ImageNet adathalmazon \cite{russakovsky2015imagenet} -- tanítják be, majd a
megtanult vizuális reprezentációkat a jóval kisebb orvosi adathalmazra finomítják
(fine-tuning). Az architektúrák fejlődése szintén hozzájárult az eredmények
javulásához: a reziduális kapcsolatokat bevezető ResNet \cite{he2016resnet}
lehetővé tette a nagyon mély hálók stabil tanítását, az EfficientNet pedig a
hálózat mélységének, szélességének és felbontásának együttes, kiegyensúlyozott
skálázásával ért el kedvező teljesítmény-paraméterszám arányt
\cite{tan2019efficientnet}.

A korai rendszerek fontos korlátja volt, hogy elsősorban dermatoszkópos, nagy
felbontású klinikai felvételeken működtek megbízhatóan, amelyek elkészítéséhez
speciális eszköz és szakértelem szükséges. Emellett a szelekciós torzítás is
jelentős problémát jelentett: az elérhető adathalmazok döntően már szakorvoshoz
került, klinikailag gyanús elváltozásokat tartalmaztak, ami tovább korlátozta a
modellek általánosíthatóságát hétköznapi felvételekre.

% =============================================================================
\section{A Vision Transformer és a figyelmi mechanizmus}
% =============================================================================

A természetesnyelv-feldolgozásban bevezetett Transformer-architektúra a
\textit{figyelmi} (attention) mechanizmusra épül, amely a bemenet elemei közötti
összes páronkénti kapcsolatot képes közvetlenül modellezni, függetlenül azok
távolságától \cite{vaswani2017attention}. Ezt az elvet a képosztályozásra a
Vision Transformer (ViT) ültette át: a képet rögzített méretű foltokra
(\textit{patch}) bontja, ezeket lineárisan beágyazza, majd egy Transformer-kódoló
dolgozza fel őket \cite{dosovitskiy2021vit}. A konvolúciós hálókkal szemben a ViT
nem rendelkezik beépített képi induktív torzítással (mint a lokalitás vagy az
eltolás-invariancia), így a globális összefüggéseket rugalmasabban tanulja meg,
cserébe viszont jellemzően nagyobb tanítóadatra vagy erőteljes előtanításra van
szüksége a jó általánosításhoz \cite{dosovitskiy2021vit}.

A ViT-architektúra a bőrgyógyászati képosztályozásban is megjelent. Dagnaw és
munkatársai például Vision Transformer és konvolúciós modelleket vetettek össze
magyarázható mesterséges intelligencia (XAI) technikákkal kiegészítve, és a
ViT-alapú, illetve a CNN-alapú modellek esetében hasonló nagyságrendű pontosságot
mértek \cite{dagnaw2024vit}. Ez a megfigyelés alátámasztja, hogy a két
architektúracsalád összevetése -- pontosan kontrollált, azonos adat- és
tanítási feltételek mellett -- önmagában is érdemi tudományos kérdés, amelyre a
dolgozat kutatási része is kitér.

% =============================================================================
\section{Az ISIC adathalmazok és versenyek}
% =============================================================================

A bőrelváltozás-osztályozás kutatását jelentősen előmozdította az annotált,
nyilvánosan elérhető adathalmazok megjelenése. A HAM10000 adathalmaz 10~015,
több forrásból származó dermatoszkópos felvételt tartalmaz a gyakori pigmentált
elváltozásokról, és egységes kiértékelési keretet biztosított a kutatóknak
\cite{tschandl2019ham10000}. A Nemzetközi Bőrképalkotási Együttműködés (ISIC)
versenysorozata tovább bővítette az elérhető képanyagot; Rotemberg és
munkatársai egy beteg-központú, a képek mellett klinikai metaadatokat is
tartalmazó adathalmazt mutattak be, hangsúlyozva a klinikai kontextus
szerepét a melanoma azonosításában \cite{rotemberg2021isic}.

A jelen dolgozat szempontjából kiemelt jelentőségű az ISIC 2024 verseny
(\textit{Skin Cancer Detection with 3D-TBP}), amely a teljes testfelszín
háromdimenziós fényképezéséből (3D total-body photography) kivágott, mintegy
400~000 elváltozás-képet tartalmaz, és binárisan -- rosszindulatú vagy
jó\-indulatú-to-intermedier kategóriába -- történő osztályozást tűz ki célul
\cite{isic2024}. Ezek a kivágott felvételek minőségükben közelebb állnak a
hétköznapi, akár okostelefonnal készíthető képekhez, mint a klasszikus,
nagy felbontású dermatoszkópos felvételek, ami a valós felhasználási
helyzethez közelíti a feladatot.

Az ISIC 2024 adathalmaz egyik meghatározó jellemzője a rendkívül erős
\textit{osztály-egyensúlytalanság}: a rosszindulatú esetek a teljes adathalmaznak
csupán töredékét (nagyságrendileg ezredrészét) teszik ki. Ennek megfelelően a
verseny nem a hagyományos pontosságot, hanem a 80\%-os érzékenység (true positive
rate) feletti tartományra korlátozott, részleges görbe alatti területet (partial
AUC) használja elsődleges metrikaként, amely orvosilag releváns módon a magas
érzékenységi tartományban méri a modellek teljesítményét \cite{isic2024}. Az
osztály-egyensúlytalanság kezelésének módszertani hátterét, valamint a
kiértékelési metrikák pontos definícióját e fejezet elméleti része tárgyalja
részletesen.

% MEGJEGYZÉS (Ervin): ide illeszthető néhány példakép az ISIC 2024 adathalmazból
% (jó- és rosszindulatú elváltozások). A versenyadat licencfeltételeit ellenőrizd
% a beillesztés előtt.
\begin{figure}[H]
  \centering
  \fbox{\begin{minipage}[c][4cm][c]{0.85\textwidth}
    \centering\itshape
    [Ábra helye] -- Példák az ISIC 2024 adathalmazból: jóindulatú (felső sor) és
    rosszindulatú (alsó sor) bőrelváltozások 3D-TBP kivágott felvételei.
  \end{minipage}}
  \caption{Reprezentatív minták az ISIC 2024 adathalmazból.}
  \label{fig:isic-samples}
\end{figure}

% =============================================================================
\section{Multimodális nagy nyelvi modellek a bőrgyógyászatban}
% =============================================================================

A legutóbbi években a multimodális nagy nyelvi modellek (multimodal large
language models, MLLM-ek) megjelenése új lehetőségeket nyitott az orvosi
képértelmezésben. Ezek a modellek -- mint a Google Gemini-családja
\cite{gemini2023} -- egyszerre képesek képi és szöveges bemenet feldolgozására,
és természetes nyelvű magyarázatot, indoklást is tudnak generálni.

A bőrgyógyászati alkalmazásukat több vizsgálat is értékelte. Cirone és
munkatársai a GPT-4~Vision és a LLaVA modelleket vetették össze melanoma és
jóindulatú elváltozások megkülönböztetésében, és azt találták, hogy a GPT-4V
lényegesen jobban teljesített (85\% pontosság a LLaVA 45\%-ával szemben),
ugyanakkor mindkét modell teljesítménye érzékelhetően romlott a sötétebb
bőrtónusok esetén -- ami a reprezentációs torzítás kockázatára és a sokszínűbb
adathalmazok szükségességére hívja fel a figyelmet \cite{cirone2024llm}. Zhou és
munkatársai célzottan bőrgyógyászati rendszert építettek (SkinGPT-4), amelyben
egy előtanított Vision Transformert egy nagy nyelvi modellel kapcsoltak össze,
és valós eseteken értékelték autonóm diagnosztikai képességét
\cite{zhou2024skingpt4}.

E munkák tanulsága kettős. Egyrészt a multimodális nyelvi modellek erőssége a
természetes nyelvű magyarázat és a felhasználóval folytatott párbeszéd, nem pedig
a megbízható, kalibrált osztályozás -- különösen erős
osztály-egyensúlytalanság mellett. Másrészt a teljesítményük populációk közötti
egyenetlensége óvatosságra int a klinikai felhasználásban. Ezek a megfontolások
közvetlenül indokolják a jelen dolgozatban követett tervezési döntést: a
\textit{klasszifikációt} egy célzottan betanított, kalibrált képi modell (ViT)
végzi, míg a nagy nyelvi modell (Gemini) szerepe a kép minőségének ellenőrzésére,
az eredmény közérthető magyarázatára és a felhasználói kérdések megválaszolására
korlátozódik. E munkamegosztás részletes tárgyalása és empirikus alátámasztása a
dolgozat kutatási, illetve mérnöki részében található.

% =============================================================================
% B) ELMÉLETI HÁTTÉR
% =============================================================================

A fejezet hátralévő része a dolgozatban felhasznált kulcstechnológiák elméleti
hátterét ismerteti: a figyelmi mechanizmustól és a Vision Transformer
felépítésétől az osztály-egyensúlytalanság kezelésén és a kiértékelési
metrikákon át a kliens--szerver kommunikáció technológiáiig.

% =============================================================================
\section{A figyelmi mechanizmus és a Transformer}
% =============================================================================

A Transformer-architektúra alapja az úgynevezett \textit{skálázott pontszorzatos
figyelem} (scaled dot-product attention), amely a bemenet elemei közötti
összefüggéseket azok távolságától függetlenül modellezi \cite{vaswani2017attention}.
A mechanizmus minden bemeneti elemhez három vektort rendel: egy lekérdezést
($Q$, query), egy kulcsot ($K$, key) és egy értéket ($V$, value). A kimenet az
értékvektorok súlyozott összege, ahol a súlyokat a lekérdezések és kulcsok
hasonlósága határozza meg:

\begin{equation}
  \mathrm{Attention}(Q, K, V) = \mathrm{softmax}\!\left(\frac{QK^{\top}}{\sqrt{d_k}}\right) V,
  \label{eq:attention}
\end{equation}

ahol $d_k$ a kulcsvektorok dimenziója; a $\sqrt{d_k}$ skálázó tényező a
pontszorzatok nagyságát stabilizálja, megakadályozva a softmax telítődését. A
gyakorlatban a modell több, párhuzamos figyelmi „fejet'' (multi-head attention)
alkalmaz, amelyek különböző reprezentációs alterekben tanulnak összefüggéseket,
így gazdagabb leírást képesek megragadni \cite{vaswani2017attention}.

A Transformer-kódoló e figyelmi rétegeket helyzetenkénti, teljesen összekötött
előrecsatolt rétegekkel (feed-forward), maradék (reziduális) kapcsolatokkal és
rétegnormalizálással kombinálja. Mivel a figyelem önmagában érzéketlen az elemek
sorrendjére, a bemenethez \textit{helyzetkódolást} (positional encoding) adnak,
amely a sorrendi információt visszacsempészi a modellbe.

% =============================================================================
\section{A Vision Transformer architektúrája}
% =============================================================================

A Vision Transformer (ViT) a természetesnyelv-feldolgozásban bevált
Transformer-kódolót ülteti át a képosztályozásra \cite{dosovitskiy2021vit}. A
feldolgozás lépései a következők:

\begin{enumerate}
  \item \textbf{Foltokra bontás és beágyazás.} A $H \times W$ méretű bemeneti
  képet rögzített méretű, nem átfedő foltokra (jellemzően $16 \times 16$ pixel)
  bontja. Minden foltot vektorrá lapít, majd egy lineáris vetítés egy rögzített
  dimenziójú beágyazási térbe képezi -- ezek a „képtokenek'', a nyelvi modellek
  szótokenjeinek analógiái.
  \item \textbf{Osztályozó token és helyzetkódolás.} A foltbeágyazások elé egy
  tanulható osztályozó tokent (\texttt{[CLS]}) fűz, amely a teljes kép összegzett
  reprezentációját hivatott hordozni. A sorrendi (térbeli) információt tanulható
  helyzetkódolás biztosítja.
  \item \textbf{Transformer-kódoló.} A tokensorozatot több, egymásra épülő
  Transformer-kódoló réteg dolgozza fel; a többfejű önfigyelem révén minden folt
  közvetlen kapcsolatba kerülhet az összes többivel, így a modell a kép globális
  összefüggéseit is megragadja.
  \item \textbf{Osztályozó fej.} A végső réteg \texttt{[CLS]} tokenjének
  reprezentációjára egy osztályozó fej (jellemzően egy- vagy kétrétegű
  perceptron) épül, amely az osztályok feletti döntést szolgáltatja.
\end{enumerate}

A konvolúciós hálókkal szemben a ViT nem rendelkezik beépített képi induktív
torzítással (lokalitás, eltolás-invariancia), ezért nagyobb adatmennyiséget vagy
erős előtanítást igényel; cserébe a globális összefüggéseket rugalmasabban
modellezi, és nagy adathalmazokon a CNN-ekkel versenyképes vagy azokat meghaladó
teljesítményt ér el \cite{dosovitskiy2021vit}.

% MEGJEGYZÉS (Ervin): ide ajánlott egy ViT-architektúra séma (foltokra bontás ->
% beágyazás -> Transformer-kódoló -> CLS -> osztályozó fej). Mivel a forráscikk
% ábrája szerzői jogvédett, érdemes saját szerkesztésű ábrát készíteni.
\begin{figure}[H]
  \centering
  \fbox{\begin{minipage}[c][4.5cm][c]{0.85\textwidth}
    \centering\itshape
    [Ábra helye] -- A Vision Transformer feldolgozási folyamata: a kép foltokra
    bontása, lineáris beágyazás, helyzetkódolás, Transformer-kódoló rétegek,
    valamint a \texttt{[CLS]} tokenre épülő osztályozó fej. (Saját szerkesztésű
    ábrát illessz be.)
  \end{minipage}}
  \caption{A Vision Transformer (ViT) sematikus felépítése.}
  \label{fig:vit-arch}
\end{figure}

% =============================================================================
\section{Átviteli tanulás és finomhangolás}
% =============================================================================

Az orvosi képadatok jellemzően korlátozott mennyiségben állnak rendelkezésre,
ami megnehezíti a nagy kapacitású modellek -- különösen a ViT -- nulláról történő
betanítását. Az \textit{átviteli tanulás} (transfer learning) erre kínál
megoldást: a modellt először egy nagy, általános képadatbázison (jellemzően az
ImageNet adathalmazon \cite{russakovsky2015imagenet}) tanítják be, ahol általános
vizuális reprezentációkat -- éleket, textúrákat, alakzatokat -- sajátít el. Az
így előtanított modellt ezután a célfeladat kisebb adathalmazán
\textit{finomhangolják} (fine-tuning).

A finomhangolás során bevett gyakorlat a \textit{differenciált tanulási ráta}
alkalmazása: az előtanított törzshálót (backbone) kis tanulási rátával frissítik,
hogy a megtanult reprezentációk ne sérüljenek, míg az újonnan inicializált
osztályozó fejet nagyobb rátával tanítják, mivel azt a feladatra nulláról kell
megtanulni. A tanulás stabilitását a \textit{bemelegítés} (warmup) is segíti,
amely a kezdeti néhány epochban fokozatosan emeli a tanulási rátát, megelőzve a
hideg fej által okozott nagy, destabilizáló lépéseket.

% =============================================================================
\section{Az osztály-egyensúlytalanság kezelése}
% =============================================================================

Az orvosi szűrési feladatok jellemzően erősen kiegyensúlyozatlanok: a pozitív
(beteg) esetek a teljes adathalmaznak csupán töredékét teszik ki. Ilyenkor a
hagyományos keresztentrópia-veszteséget a sok, könnyen osztályozható negatív
példa uralja, és a modell hajlamos a többségi osztály felé torzítani. A
szakirodalom több, egymást kiegészítő technikát kínál e probléma kezelésére.

\paragraph{Súlyozott mintavételezés.} Az egyik megközelítés a tanító batch-ek
osztályeloszlásának kiegyensúlyozása: minden mintát az osztálya gyakoriságával
fordítottan arányos súllyal választunk, így a ritka pozitív esetek a tanítás
során felülreprezentálttá válnak.

\paragraph{Fókuszveszteség (focal loss).} A Lin és munkatársai által bevezetett
fókuszveszteség a keresztentrópiát úgy alakítja át, hogy a jól osztályozott
példák hozzájárulását lecsökkenti, és a tanulást a nehéz, félreosztályozott
esetekre összpontosítja \cite{lin2017focal}:

\begin{equation}
  \mathrm{FL}(p_t) = -\,\alpha_t \,(1 - p_t)^{\gamma}\, \log(p_t),
  \label{eq:focal}
\end{equation}

ahol $p_t$ a helyes osztályhoz rendelt becsült valószínűség, $\gamma \ge 0$ a
fókuszáló paraméter (amely a könnyű példák súlyát csökkenti), $\alpha_t$ pedig
egy osztálysúly, amely a pozitív és negatív esetek közötti egyensúlyt
finomhangolja. A $\gamma = 0$ választás visszaadja a súlyozott
keresztentrópiát.

\paragraph{Kétfázisú tanítás.} A kiegyensúlyozott mintavételezés mellett betanított
modell kimeneti valószínűségei eltérhetnek a valós, erősen kiegyensúlyozatlan
előfordulási aránytól. Ezt egy második, kalibráló fázis korrigálhatja, amelyben
a modellt a valós eloszlású adaton, kisebb tanulási rátával tanítjuk tovább, így
a kimeneti valószínűségek a tényleges priorhoz igazodnak.

% =============================================================================
\section{Kiértékelési metrikák}
% =============================================================================

Erősen kiegyensúlyozatlan feladatban a nyers pontosság (accuracy) félrevezető,
hiszen a többségi osztály triviális megjósolásával is magas érték érhető el.
Ezért orvosi szűrésnél a tévesztés-mátrixból (confusion matrix) származtatott,
osztályérzékeny metrikákat alkalmazunk. Jelölje $TP$, $TN$, $FP$ és $FN$ rendre a
valódi pozitív, valódi negatív, hamis pozitív és hamis negatív esetek számát.

\begin{align}
  \text{Érzékenység (sensitivity)} &= \frac{TP}{TP + FN}, &
  \text{Specificitás (specificity)} &= \frac{TN}{TN + FP}, \label{eq:sens-spec}\\
  \text{PPV} &= \frac{TP}{TP + FP}, &
  \text{NPV} &= \frac{TN}{TN + FN}. \label{eq:ppv-npv}
\end{align}

Az \textbf{érzékenység} azt méri, hogy a modell a valóban rosszindulatú eseteknek
mekkora hányadát ismeri fel -- szűrőeszköznél ez a legkritikusabb mutató, mivel a
fel nem ismert rák (hamis negatív) a legsúlyosabb kimenetel. A \textbf{negatív
prediktív érték} (NPV) azt fejezi ki, hogy a „jóindulatú'' besorolás mennyire
megbízható.

\paragraph{ROC-görbe és AUC.} A bináris osztályozók döntési küszöbtől független
értékelésére a \textit{működési jelleggörbét} (Receiver Operating Characteristic,
ROC) használjuk, amely a valódi pozitív arányt (érzékenység) a hamis pozitív
arány ($FPR = FP/(FP+TN)$) függvényében ábrázolja a küszöb változtatásával. A
görbe alatti terület (Area Under the Curve, AUC) egyetlen számba sűríti a
rangsorolási képességet: értéke 0,5 (véletlen találgatás) és 1,0 (tökéletes
elkülönítés) között mozog.

\paragraph{Részleges AUC (pAUC).} Erős egyensúlytalanságnál és klinikai
kontextusban gyakran csak a görbe egy meghatározott szakasza érdekes. Az ISIC
2024 verseny elsődleges metrikája a 80\%-os érzékenység feletti tartományra
korlátozott részleges AUC (partial AUC), amely kifejezetten a klinikailag
elvárt, magas érzékenységi tartományban méri a teljesítményt; a metrika
értéke $[0,\ 0{,}2]$ tartományba esik, ahol $0{,}2$ a tökéletes osztályozó
értéke \cite{isic2024}.

\paragraph{Kalibráció és Brier-pontszám.} A jó rangsorolás (magas AUC) nem
garantálja, hogy a modell által kiadott valószínűségek megbízhatóak. A
\textit{kalibráció} azt méri, hogy a becsült valószínűségek mennyire felelnek
meg a tényleges bekövetkezési aránynak; Guo és munkatársai rámutattak, hogy a
modern neurális hálók gyakran rosszul kalibráltak \cite{guo2017calibration}. A
kalibráció egyik összegző mérőszáma a Brier-pontszám, a becsült valószínűség és a
tényleges címke közötti átlagos négyzetes eltérés:

\begin{equation}
  \mathrm{Brier} = \frac{1}{N} \sum_{i=1}^{N} (p_i - y_i)^2,
  \label{eq:brier}
\end{equation}

ahol $p_i$ a $i$-edik mintára becsült valószínűség, $y_i \in \{0,1\}$ a valódi
címke; az alacsonyabb érték jobb kalibrációt jelez.

% =============================================================================
\section{Statisztikai elemzés}
% =============================================================================

Mivel a teszthalmazban a pozitív esetek száma korlátozott, a pontbecslések
(érzékenység, specificitás, AUC) mellett azok bizonytalanságát is jelölni kell.
Az arányjellegű metrikákhoz (érzékenység, specificitás) az egzakt, binomiális
eloszláson alapuló Clopper--Pearson-féle konfidencia-intervallum ad megbízható
becslést kis mintaszám esetén is. Az AUC esetében a DeLong-féle, nemparaméteres
módszer szolgáltat konfidencia-intervallumot, illetve teszi lehetővé két modell
AUC-értékének páros összehasonlítását; a módszer egy gyors, korszerű
implementációját Sun és Xu adta meg \cite{sun2014delong}. E statisztikai eszközök
biztosítják, hogy a modellek közötti különbségekről tett állítások ne csupán a
mintavételi zaj következményei legyenek.

% =============================================================================
\section{Kommunikációs technológiák: SSE és az MCP}
% =============================================================================

A dolgozatban bemutatott rendszer valós idejű, token-szintű válaszközvetítést és
egy nyelvi modellt kiszolgáló komponensekkel való kommunikációt igényel; ezt két
technológia teszi lehetővé.

\paragraph{Szerver-küldött események (SSE).} A szerver-küldött események
(Server-Sent Events) a HTML szabvány részét képező, egyirányú, szerverről
kliensre irányuló adatfolyam-mechanizmus egyetlen, tartós HTTP-kapcsolaton
keresztül \cite{whatwg_sse}. Az SSE különösen alkalmas a nagy nyelvi modellek
token-szintű, fokozatos válaszainak továbbítására, mivel a kliens a szöveget
generálással egyidejűleg, darabonként jeleníti meg, jelentősen javítva a
felhasználói élményt.

\paragraph{Model Context Protocol (MCP).} A Model Context Protocol az Anthropic
által 2024-ben közzétett nyílt szabvány, amely egységesíti a nagy nyelvi
modellek és a külső eszközök, adatforrások közötti kommunikációt \cite{mcp2024}.
Az MCP a JSON-RPC 2.0 üzenetformátumra épül \cite{jsonrpc2013}, és többféle
átviteli réteget támogat, többek között a helyi folyamatok közötti szabványos
be- és kimeneti (stdio) csatornát. A dolgozat rendszerében az MCP biztosítja a
backend és a modellt kiszolgáló komponens közötti, jól strukturált
kommunikációt; e megoldás részletes tárgyalása a mérnöki részben található.

